\ysrc{158}
然后我们如下定义与 $\mathcal{D}$ 相伴的游历算子
\[
\mathrm{Id}_{\mathcal{C}} = T_{\mathbf{1}} \xrightarrow{T_\alpha} T_V \xrightarrow{\gamma} T_V \xrightarrow{T_\beta} T_{\mathbf{1}} = \mathrm{Id}_{\mathcal{C}},
\]
其中对 $T_\alpha$（相应地 $T_\beta$）我们用了函子性 $\varnothing\to I$（相应地 $I\to\varnothing$）。让 $\gamma$ 变化，四元组 $(I,V,\alpha,\beta)$ 给出一个映射
\[
S(V,\alpha,\beta) : W^I \to \End(\mathrm{Id}).
\]
注意，若 $g\in(\hatG\rtimes Q)^I$，则它定义了一个标量
\[
\mathbf{1}\xrightarrow{\ \alpha\ }V\xrightarrow{\ g\ }V\xrightarrow{\ \beta\ }\mathbf{1}.
\]
让 $g$ 变化，这给出了 $\mathcal{O}(\hatG\backslash(\hatG\rtimes Q)^I/\hatG)$ 中的一个元素 $f(I,V,\alpha,\beta)$。定理 5.9 中的态射于是由
\[
f(V,\alpha,\beta)\mapsto S(V,\alpha,\beta)
\]
并验证各种相容性而构造出来。

\subsubsection{L-参数}
于是我们如下构造态射（取 $\Lambda=\Q_\ell[\sqrt q]$；当 $\ell$ 对 $\hatG$ 非常好时，取 $\Lambda=\Z_\ell[\sqrt q]$）：
\[
\underbrace{\mathfrak Z^{\mathrm{spec}}(G,\Lambda)}_{\text{谱稳定中心}}
\longrightarrow
\underbrace{\mathfrak Z^{\mathrm{geo}}(G,\Lambda)}_{\text{几何稳定中心}}
\longrightarrow
\underbrace{\mathfrak Z(G(E),\Lambda)}_{\text{Bernstein 中心}}.
\]
若 $\ell$ 对 $\hatG$ 不是非常好的素数，则它只在泛同胚的意义下才有定义。这是 Helm 与 Moss 关于族中局部 Langlands 对应的工作 [12] 的推广。利用上述中心之间的态射，我们证明了下述定理。

\begin{thm}{5.10}\label{thm:a3-5-10}
设 $L$ 为 $\overline{\Q}_\ell$ 或 $\overline{\F}_\ell$ 之一。设 $\pi$ 是 $G(E)$ 的以 $L$ 为系数的光滑不可约表示。我们可以构造其半单 Langlands 参数 $\varphi_\pi$。此外它还满足：
\begin{enumerate}[label=(\arabic*),leftmargin=2.6em,itemsep=0.2em]
\item 它与抛物诱导相容。
\item 它与标量的 Weil 限制相容。\tn{底本此处作「Is is compatible」，疑为「It is compatible」之误。}
\item 它与乘积相容，即 $\varphi_{\pi_1\boxtimes\pi_2}=\varphi_{\pi_1}\times\varphi_{\pi_2}$。
\item 若 $G$ 是环面，则它由局部类域论给出。\tn{底本此处作「local class field」，疑脱 theory 一词。}
\item 它与 $\mathrm{GL}_n$ 的半单局部 Langlands 对应一致（[11]、[13]、[27]、[24]）。
\end{enumerate}
\end{thm}

第 (5) 条是利用与 Lubin--Tate 空间上同调的相容性来验证的。与 Harris 与 Taylor 的工作 [11] 的联系正是在此处建立的。此外还需指出，关于 $\mathrm{Gsp}_4$ 的 Gan--Takeda 局部 Langlands 对应的相容性已在 [10] 中验证。

\subsubsection{对 $\ell$ 的无关性}
下述猜想是自然的。注意，$\mathfrak Z^{\mathrm{spec}}(G)$（乃至 $\LocSys_{\hatG}$）与通常的 Bernstein 中心 $\mathfrak Z(G(E))$ 都自然地定义为平坦的 $\Z[\frac1p]$-代数。

\begin{conj}{5.11}\label{conj:a3-5-11}
设 $N$ 为 $p$ 与那些对 $\hatG$ 不是非常好素数的 $\ell$ 之积。存在 $\Z[\frac1N]$-代数的态射
\[
\mathfrak Z^{\mathrm{spec}}(G,\Z[\tfrac1N])\longrightarrow \mathfrak Z(G(E),\Z[\tfrac1N]),
\]
它对所有 $\ell\neq p$ 且对 $\hatG$ 非常好的素数都诱导出上述中心之间的态射。
\end{conj}

\ysrc{159}
\subsection{谱作用}
设 $\Lambda\in\{\Z_\ell[\sqrt q],\Q_\ell[\sqrt q]\}$。若 $\Lambda=\Z_\ell[\sqrt q]$，则假定 $\ell$ 是非常好的素数。下面我们说明如何把态射
\[
\mathfrak Z^{\mathrm{spec}}(G,\Lambda)\longrightarrow\mathfrak Z^{\mathrm{geo}}(G,\Lambda)
\]
的构造升级，正是它使我们能够构造半单 $L$-参数。这将把参数的自同构纳入考虑，为此我们将在更高范畴的框架中工作。

\subsubsection{模表示论}
（[9, 第 VIII.5 章]）这里的主要结果如下。此处我们在 $\Z_\ell$ 上工作。

\begin{thm}{5.12}\label{thm:a3-5-12}
设 $\ell$ 对 $\hatG$ 是非常好的素数。则态射
\[
\varinjlim_{(n,F_n)\in\mathfrak{F}}\mathcal{O}(Z^1(F_n,\hatG))\longrightarrow\mathcal{O}(Z^1(W_E,\hatG))
\]
是可表现稳定 $\infty$-范畴 $\mathrm{Ind}\,\Perf(B\hatG)$ 中的同构。
\end{thm}

这里 $\Perf(B\hatG)$ 是代数叠 $B\hatG$ 上完美复形的 $\infty$-范畴。其同伦范畴就是 $\hatG$ 在有限自由 $\Z_\ell$-模上的代数表示的有界复形所成的范畴。$\mathcal{O}(Z^1(F_n,\hatG))$ 与 $\mathcal{O}(Z^1(W_E,\hatG))$ 这两个对象都可以典范地看作 Ind-完美复形：把它们写成其有限型的、$\hatG$-稳定的 $\Z_\ell$-子模的归纳极限即可。该定理说的是，当把 $\varinjlim_{\mathfrak{F}}$ 作用到 $(\mathcal{O}(F_n,\hatG))_{(n,F_n)\in\mathfrak{F}}$ 上时「没有高阶导出函子」。\tn{底本此处印作「$((\mathcal{O}(F_n,\hatG))_{(n,F_n)\in\mathfrak{F}}$」，括号不配对，疑为多一左括号。}它立刻蕴含定理 5.8。事实上，这表明对一切 $i\ge 0$ 都有
\[
\varinjlim_{(n,F_n)\in\mathfrak{F}}H^i\bigl(\hatG,\mathcal{O}(Z^1(F_n,\hatG))\bigr)\xrightarrow{\ \sim\ }H^i\bigl(\hatG,\mathcal{O}(Z^1(W_E,\hatG))\bigr).
\]
由此我们可以证明，对 $i>0$ 也有
\[
H^i\bigl(\hatG,\mathcal{O}(Z^1(W_E,\hatG))\bigr)=0,
\]
而关于 $\mathcal{O}(Z^1(F_n,\hatG))$ 的结论则可容易地由模表示论中一些已知结果推出。

上述结果还可以按下述方式得到进一步加强。

\begin{thm}{5.13}\label{thm:a3-5-13}
假定或者我们在 $\Q_\ell$ 上工作，或者我们在 $\Z_\ell$ 上工作且 $\ell$ 对 $\hatG$ 是非常好的素数。则 $\infty$-范畴 $\Perf(\LocSys_{\hatG})$ 由 $\Perf(B\hatG)$ 在锥与收缩元下生成。
\end{thm}

这正是我们用来定义谱作用的东西。

\subsubsection{谱作用}
利用更高范畴的方法、定理 5.12 与定理 5.13，我们可以按下述方式加强定理 5.9。设 $\Lambda$ 是 $\Q_\ell[\sqrt q]$ 的某个有限次扩张中的整数环。我们固定 $W_E$ 的一个有限商，使得对 $\hatG$ 的作用经过它分解。

\begin{thm}{5.14}\label{thm:a3-5-14}
设 $\ell$ 对 $\hatG$ 是非常好的素数。设 $\mathcal{C}$ 是小幂等完备的 $\Lambda$-线性 $\infty$-范畴。则下述数据是等价的：
\begin{enumerate}[label=(\arabic*),leftmargin=2.6em,itemsep=0.2em]
\item 对有限集 $I$ 函子性地给出一个正合的 $\Rep_\Lambda(Q^I)$-线性幺半函子
\[
\Rep_\Lambda((\hatG\rtimes Q)^I)\longrightarrow \End_\Lambda(\mathcal{C})^{BW_E^I}.
\]
\ysrc{160}
\item $\Perf(\LocSys_{\hatG})$ 在 $\mathcal{C}$ 上的一个 $\Lambda$-线性作用，使得对每个 $X\in\mathcal{C}$，其在 $X$ 上的相应作用经过 $\Perf(Z^1(W_E/P,\hatG)_\Lambda/\hatG_\Lambda)$ 分解，其中 $P$ 是 $W_E$ 的野惯性的某个开子群。
\end{enumerate}
\end{thm}

同一断言在 $\Q_\ell[\sqrt p]$ 的有限次扩张上也成立，且对 $\ell$ 无任何限制。\tn{底本此处作 $\Q_\ell[\sqrt p]$，与本定理前半段的 $\Q_\ell[\sqrt q]$ 不一致，疑为笔误。}

在上述命题中，$[Z^1(W_E/P,\hatG)_\Lambda/\hatG_\Lambda]$ 是 $\LocSys_{\hatG}$ 的一个开闭、拟紧子叠。对于定理中关于 $X$ 上作用的那一相应条件，我们使用「紧支撑作用」这一术语。

在此定理中，从 (2) 到 (1) 的对应由（显然的）$\Rep(Q^I)$-线性幺半函子
\[
\Rep_\Lambda(\hatG\rtimes Q)^I\longrightarrow \Perf\bigl(Z^1(W_E,\hatG)^I_\Lambda/\hatG_\Lambda\bigr)^{BW_E^I}
\]
与谱作用复合给出。

定理 5.7 中出现的一切对象都有自然的 $\infty$-范畴升级。特别地，我们定义一个稳定 $\Lambda$-线性 $\infty$-范畴
\[
\mathcal{D}_{lis}(\Bun_G,\Lambda),
\]
其同伦范畴是 $D_{lis}(\Bun_G,\Lambda)$。这使我们能够证明下述定理。

\begin{thm}{5.15 (谱作用)}\label{thm:a3-5-15}
假定或者 $\Lambda$ 是 $\Q_\ell[\sqrt q]$ 的一个扩张，或者 $\Lambda$ 是 $\Q_\ell[\sqrt q]$ 的某个有限次扩张中的整数环（在后一情形假定 $\ell$ 对 $\hatG$ 非常好）。则存在 $\Perf(\LocSys_{\hatG}/\Lambda)$ 在 $\mathcal{D}_{lis}(\Bun_G,\Lambda)^\omega$ 上的一个自然的 $\Lambda$-线性紧支撑作用。
\end{thm}

与前面一样，该作用由其与 Hecke 算子作用的相容性唯一地确定。若忘掉 $W_E$-作用，则与 $V\in\Rep_\Lambda(\hatG)$ 相伴的 Hecke 算子 $T_V$ 由态射
\[
\LocSys_{\hatG}\longrightarrow B\hatG
\]
推出，而后者诱导出幺半函子
\[
\Rep_\Lambda(\hatG)\longrightarrow \Perf(\LocSys_{\hatG}).
\]

\subsubsection{对尖点参数的应用}
我们给出谱作用的一个应用实例。我们把自己置于 $\overline{\Q}_\ell$ 上。按定义，一个 $L$-参数 $\varphi:W_E\to{}^LG(\overline{\Q}_\ell)$ 是尖点的（cuspidal），若它是半单的且 $S_\varphi/Z(\hatG)^\Gamma$ 有限。容易验证这定义了一个连通分支
\[
\mathcal{C}_\varphi\subset\LocSys_{\hatG},
\]
即 $\varphi$ 的非分歧扭化所成的开闭子叠。存在态射
\[
[\Spec(\overline{\Q}_\ell)/S_\varphi]\longrightarrow\mathcal{C}_\varphi\subset\LocSys_{\hatG},
\]
它是闭浸入。态射 $\mathfrak Z^{\mathrm{spec}}\to\mathfrak Z^{\mathrm{geo}}$ 把与 $\mathcal{C}_\varphi$ 相伴的幂等元送到 $\mathfrak Z^{\mathrm{geo}}$ 中的一个幂等元。这定义了一个直和项
\[
\mathcal{D}_{lis}(\Bun_G,\overline{\Q}_\ell)^\omega[\varphi]\subset\mathcal{D}_{lis}(\Bun_G,\overline{\Q}_\ell)^\omega.
\]
我们来分析这一点。对任意 Schur 不可约的 $A\in D_{lis}(\Bun_G,\overline{\Q}_\ell)[\varphi]$（即 $\End(A)=\overline{\Q}_\ell$），游历算子在 $A$ 上按标量作用，而这些标量由 $\varphi$ 的某个非分歧扭化所确定。

\ysrc{161}
对一切 $[b]\in B(G)$，它们在 $(i^b)^*A$ 上通过游历代数的同一个特征标作用。由 Langlands 参数的构造与抛物诱导的相容性，再结合 $\varphi$ 的尖点性（它不经过 $^LG$ 的任何抛物子群分解），我们推出：若 $b$ 不是基本的，则 $(i^b)^*A=0$。由这一论证我们立即得到下述结果。

\begin{thm}{5.16 (尖点参数的洁净性)}\label{thm:a3-5-16}
设 $\pi$ 是 $G(E)$ 的不可约表示，且 $\varphi$ 是尖点的。则 $\pi$ 是超尖点的，并且
\[
(i^1)_!\mathcal{F}_\pi=R(i^1)_*\mathcal{F}_\pi.
\]
\end{thm}

为简化起见，现在设 $Z(\hatG)^\Gamma$ 有限。由上述论证我们可以推出，经由 $(i^1)_!$，
\[
D_{lis}(\Bun_G,\overline{\Q}_\ell)[\varphi]=\bigoplus_{[b]\ \text{basic}}\ \bigoplus_{\substack{\pi\ \text{supercusp. of }G_b(E)\\ \varphi_\pi=\varphi}}\Perf(\overline{\Q}_\ell)\otimes\pi.
\]
现在我们使用谱作用：存在 $\Rep_{\overline{\Q}_\ell}(S_\varphi)$ 在 $D_{lis}(\Bun_G,\overline{\Q}_\ell)[\varphi]$ 上的幺半作用。对任意 $\rho\in\mathrm{Irr}_{\overline{\Q}_\ell}(S_\varphi)$ 有
\[
\rho|_{Z(\hatG)^\Gamma}\in X^*(Z(\hatG)^\Gamma)=\pi_1(G)_\Gamma.
\]
按谱作用的构造，对基本的 $[b]$，若 $[b']$ 也是基本的且
\[
\kappa(b')=\kappa(b)+\rho|_{Z(\hatG)^\Gamma},
\]
则有
\[
\rho*(-):\bigoplus_{\substack{\pi\ \text{supercusp. of }G_b(E)\\ \varphi_\pi=\varphi}}\Perf(\overline{\Q}_\ell)\otimes\pi\longrightarrow\bigoplus_{\substack{\pi\ \text{supercusp. of }G_{b'}(E)\\ \varphi_\pi=\varphi}}\Perf(\overline{\Q}_\ell)\otimes\pi.
\]
这种平移是 Jacquet--Langlands 对应的一种形式。受此启发，我们提出下述猜想。

\begin{conj}{5.17}\label{conj:a3-5-17}
设 $G$ 拟分裂，并固定一个 Whittaker 数据 $(B,\psi)$。
\begin{enumerate}[label=(\arabic*),leftmargin=2.6em,itemsep=0.2em]
\item 存在唯一的 $G(E)$ 的不可约表示 $\pi$，其参数为 $\varphi$，且关于 $(B,\psi)$ 是泛型的（generic）。
\item 幺半作用
\[
\Perf(\Rep_{\overline{\Q}_\ell}(S_\varphi))=\bigoplus_{n\in\Z}\Rep_{\overline{\Q}_\ell}(S_\varphi)[n]\longrightarrow D_{lis}(\Bun_G,\overline{\Q}_\ell)^\omega[\varphi],\qquad \rho[n]\longmapsto\rho*\pi[n]
\]
是等价。
\end{enumerate}
\end{conj}

在下一节中，我们把这一猜想从 $D_{lis}(\Bun_G,\overline{\Q}_\ell)^\omega[\varphi]$（$\varphi$ 尖点）推广到整个范畴 $D_{lis}(\Bun_G,\overline{\Q}_\ell)^\omega$。

\ysrc{162}
\section{范畴几何化猜想}
本节我们解释 [9] 的主要猜想，并给出它的一些推论。

\subsection{Arinkin--Gaitsgory 奇异支撑条件}
（[9, 第 VIII.2.2 章]）设 $\mathcal{X}\to S$ 是正则概形 $S$ 上局部完全交的代数叠。我们可以看它的余切复形
\[
\mathbb{L}_{\mathcal{X}/S}\in\Perf^{[-1,1]}(\mathcal{O}_{\mathcal{X}}).
\]
Arinkin 与 Gaitsgory（[2]）定义了奇点叠
\[
\Sing_{\mathcal{X}/S}:=\mathcal{H}^{-1}(\mathbb{L}_{\mathcal{X}/S})\longrightarrow\mathcal{X}.
\]
这是 $\mathcal{X}$ 上带有 $\mathbb{G}_m$ 作用的交换群概形。对 $\mathcal{E}\in D^b_{coh}(\mathcal{O}_{\mathcal{X}})$，他们定义其奇异支撑
\[
\mathrm{SingSupp}(\mathcal{E})\subset\Sing_{\mathcal{X}/S},
\]
即奇点叠的一个 Zariski 闭、$\mathbb{G}_m$-不变（亦即锥形）的子集。这是凝聚语境中某种「微支撑」（microsupport）。它在 $\mathcal{X}$ 中的像就是 $\mathcal{E}$ 的支撑。他们证明了下述结果。

\begin{thm}{6.1 (Arinkin--Gaitsgory)}\label{thm:a3-6-1}
奇异支撑 $\mathrm{SingSupp}(\mathcal{E})$ 包含于零截面 $\{0\}\subset\Sing_{\mathcal{X}/S}$ 之中，当且仅当 $\mathcal{E}$ 是完美复形。
\end{thm}

这是下述事实的凝聚类比：一个反常层的特征循环包含于零截面之中，当且仅当它是局部系统。

现在，见 §5.1.3，我们有
\[
\Sing_{\LocSys_{\hatG}/\Z_\ell}\subset[\frakg^*/\hatG]\times_{B\hatG}\LocSys_{\hatG}.
\]
我们定义
\[
\Coh_{Nilp}(\LocSys_{\hatG})\subset D^b_{coh}(\LocSys_{\hatG},\mathcal{O})
\]
为满足下述条件的复形所成的子范畴：其支撑拟紧，即支撑在 $\LocSys_{\hatG}$ 的有限个连通分支上，且其奇异支撑包含于幂零锥
\[
[\mathcal{N}^*_{\hatG}/\hatG]\times_{B\hatG}\LocSys_{\hatG}
\]
之中。注意，在把 $\ell$ 可逆化之后，这一条件自动成立。

\subsection{该猜想}
（[9, 第 X.3 章]）设 $G$ 拟分裂。设 $U$ 是 $G$ 的某个 Borel 子群的幂幺根，而
\[
\psi:U(E)\to\overline{\Z}_\ell^\times
\]
是非退化特征标。考虑 Whittaker 层
\[
\mathcal{W}_\psi=(i^1)_!\bigl(c\text{-}\Ind_{U(E)}^{G(E)}\psi\bigr).
\]
它不是 $D_{lis}(\Bun_G,\overline{\Z}_\ell)$ 的紧对象。尽管如此，对 $\mathcal{F}\in\Perf(\LocSys_{\hatG})$，我们仍可定义 $\mathcal{F}$ 与此对象的谱作用 $\mathcal{F}*\mathcal{W}_\psi\in D_{lis}(\Bun_G,\overline{\Z}_\ell)$：把 $c\text{-}\Ind_{U(E)}^{G(E)}\psi$ 写成有限型表示的余极限即可。

\ysrc{163}
下述猜想是猜想 5.17 的升级，而后者在某种意义上正是它的「玩具模型」。我们在整系数下工作，因此假定 $\ell$ 对 $\hatG$ 是非常好的素数。在 $\overline{\Q}_\ell$ 上我们提出同一猜想，而不对素数 $\ell$ 作此限制。

\begin{conj}{6.2 (范畴几何化猜想)}\label{conj:a3-6-2}
函子
\[
\Perf(\LocSys_{\hatG})\longrightarrow D_{lis}(\Bun_G,\overline{\Z}_\ell),\qquad \mathcal{F}\longmapsto\mathcal{F}*\mathcal{W}_\psi
\]
限制到具有拟紧支撑的完美复形上时取值于紧对象，并且延拓为等价
\[
\Coh_{Nilp}(\LocSys_{\hatG})\xrightarrow{\ \sim\ }D_{lis}(\Bun_G,\overline{\Z}_\ell)^\omega,
\]
它与谱作用相容。
\end{conj}

这就是我们所设想的局部 Langlands 对应的最终形式。

\subsection{若干推论}
下面我们指出范畴几何化猜想的一些推论。

\subsubsection{稳定中心之间的等同}
设 $\Lambda\in\{\Z_\ell[\sqrt q],\Q_\ell[\sqrt q]\}$。在整系数情形我们进一步假定 $\ell$ 对 $\hatG$ 是非常好的素数。几何化猜想的全忠实性部分蕴含上述复合是同构。人们还可以希望用稳定分布来描述这个中心，但这与我们的工作并无关联。

\subsubsection{函子性的核}
（[9, 第 X.1 章]）此处我们在 $\overline{\Q}_\ell$ 上工作。设 $H$ 与 $G$ 是 $E$ 上的拟分裂约化群。我们为这两个群都固定 Whittaker 数据。设给定一个 $L$-态射
\[
f:{}^LH\longrightarrow{}^LG.
\]
范畴猜想蕴含：谱一侧由态射 $f_*:\LocSys_{\widehat{H}}\to\LocSys_{\widehat{G}}$ 给出的函子性，将在几何一侧给出一个函子
\[
D_{lis}(\Bun_H,\overline{\Q}_\ell)\longrightarrow D_{lis}(\Bun_G,\overline{\Q}_\ell).
\]
我们证明这样的函子自动由一个核给出
\[
A_f\in D_{lis}(\Bun_H\times\Bun_G,\overline{\Q}_\ell),
\]
即一个函子性的核。再投影到「经典表示论的部分」，我们还得到一个函子 $D(H(E),\overline{\Q}_\ell)\to D(G(E),\overline{\Q}_\ell)$。

\begin{thm}{6.3}\label{thm:a3-6-3}
范畴几何化猜想蕴含拟分裂约化群的局部 Langlands 函子性，即与 $L$-态射 ${}^LH\to{}^LG$ 相伴的函子 $D(H(E),\overline{\Q}_\ell)\to D(G(E),\overline{\Q}_\ell)$。它由与该 $L$-态射相伴的一个「自然的」核 $A\in D_{lis}(\Bun_H\times\Bun_G,\overline{\Q}_\ell)$ 给出。
\end{thm}

\ysrc{164}
\section{一些最后的想法}
归根结底，局部 Langlands 纲领所适用的自然对象看来并不是 $G(E)$ 的光滑表示，而是 $\mathcal{D}_{lis}(\Bun_G,\overline{\Q}_\ell)$ 的对象。典型地，对 Schur 不可约的 $A\in\mathcal{D}_{lis}(\Bun_G,\overline{\Q}_\ell)$，我们可以附上其半单 Langlands 参数
\[
\varphi_A:W_E\longrightarrow{}^LG(\overline{\Q}_\ell).
\]
此外，见 §3.4，有限型、容许、Zelevinsky 对合这些概念都自然地推广为 $\mathcal{D}_{lis}(\Bun_G,\overline{\Q}_\ell)$ 中的几何概念。如前所见，局部 Langlands 函子性在 $\mathcal{D}_{lis}(\Bun_G,\overline{\Q}_\ell)$ 的层面上自然地由核定义。

这就引出下述问题：自守表示是整体 Langlands 纲领所适用的自然对象吗？正如我们在局部情形已经看到的，从表示论的观点看，自然对象并不是 $G(E)$ 的表示，而是所有 $G_b(E)$（$[b]\in B(G)$）的表示同时在一起。整体 Kottwitz 集是存在的（[15]），因此自然要问：我们是否应当同时考虑所有相应的 $G_b$ 的自守表示？

从 $\mathcal{D}_{lis}(\Bun_G,\overline{\Q}_\ell)$ 的观点看，整体曲线并不存在，情况更为神秘。不过仍要指出，前述工作的阿基米德类比尚待寻找。

\section*{参考文献}
\begin{refs}
\item[{[1]}] J. Ansch\"utz, \emph{Reductive group schemes over the Fargues-Fontaine curve}, Math. Ann. \textbf{374} (2019), no. 3-4, 1219--1260, DOI 10.1007/s00208-018-1785-6. MR3985110
\item[{[2]}] D. Arinkin and D. Gaitsgory, \emph{Singular support of coherent sheaves and the geometric Langlands conjecture}, Selecta Math. (N.S.) \textbf{21} (2015), no. 1, 1--199, DOI 10.1007/s00029-014-0167-5. MR3300415
\item[{[3]}] B. Bhatt and P. Scholze, \emph{Projectivity of the Witt vector affine Grassmannian}, Invent. Math. \textbf{209} (2017), no. 2, 329--423, DOI 10.1007/s00222-016-0710-4. MR3674218
\item[{[4]}] C. Birkbeck, T. Feng, D. Hansen, S. Hong, Q. Li, A. Wang, and L. Ye, \emph{Extensions of vector bundles on the Fargues-Fontaine curve}, J. Inst. Math. Jussieu \textbf{21} (2022), no. 2, 487--532, DOI 10.1017/S1474748020000183. MR4386821
\item[{[5]}] J.-F. Dat, D. Helm, R. Kurinczuk, and G. Moss, \emph{Moduli of Langlands parameters}, J. Eur. Math. Soc. (JEMS) \textbf{27} (2025), no. 5, 1827--1927, DOI 10.4171/jems/1599. MR4889237
\item[{[6]}] L. Fargues, \emph{G-torseurs en th\'eorie de Hodge p-adique} (French, with English and French summaries), Compos. Math. \textbf{156} (2020), no. 10, 2076--2110, DOI 10.1112/s0010437x20007423. MR4179595
\item[{[7]}] L. Fargues, \emph{Simple connexit\'e des fibres d'Abel-Jacobi et corps de classes local} (French, with English and French summaries), Ann. Sci. \'Ec. Norm. Sup\'er. (4) \textbf{53} (2020), no. 1, 89--124, DOI 10.24033/asens.2418. MR4093441
\item[{[8]}] L. Fargues and J.-M. Fontaine, \emph{Courbes et fibr\'es vectoriels en th\'eorie de Hodge p-adique} (French, with English and French summaries), with a preface by Pierre Colmez, Ast\'erisque \textbf{406} (2018), xiii+382. MR3917141
\item[{[9]}] L. Fargues and P. Scholze, \emph{Geometrization of the local Langlands correspondence}, arXiv:2102.13459, 2021, to appear in Ast\'erisque.
\item[{[10]}] L. Hamann, \emph{Compatibility of the Fargues--Scholze and Gan--Takeda local Langlands}, Compos. Math. \textbf{161} (2025), no. 9, 2202--2271, DOI 10.1112/S0010437X25102455. MR4978375
\item[{[11]}] M. Harris and R. Taylor, \emph{The geometry and cohomology of some simple Shimura varieties}, with an appendix by Vladimir G. Berkovich, Annals of Mathematics Studies, vol. 151, Princeton University Press, Princeton, NJ, 2001. MR1876802
\item[{[12]}] D. Helm and G. Moss, \emph{Converse theorems and the local Langlands correspondence in families}, Invent. Math. \textbf{214} (2018), no. 2, 999--1022, DOI 10.1007/s00222-018-0816-y. MR3867634
\ysrc{165}
\item[{[13]}] G. Henniart, \emph{Une preuve simple des conjectures de Langlands pour GL(n) sur un corps p-adique} (French, with English summary), Invent. Math. \textbf{139} (2000), no. 2, 439--455, DOI 10.1007/s002220050012. MR1738446
\item[{[14]}] K. S. Kedlaya and R. Liu, \emph{Relative p-adic Hodge theory: foundations} (English, with English and French summaries), Ast\'erisque \textbf{371} (2015), 239. MR3379653
\item[{[15]}] R. Kottwitz, \emph{$B(G)$ for all local and global fields}, arXiv:1401.5728 [math.RT].
\item[{[16]}] R. E. Kottwitz, \emph{Isocrystals with additional structure}, Compositio Math. \textbf{56} (1985), no. 2, 201--220. MR809866
\item[{[17]}] V. Lafforgue, \emph{Chtoucas pour les groupes r\'eductifs et param\'etrisation de Langlands globale} (French), J. Amer. Math. Soc. \textbf{31} (2018), no. 3, 719--891, DOI 10.1090/jams/897. MR3787407
\item[{[18]}] A.-C. Le Bras, \emph{Espaces de Banach-Colmez et faisceaux coh\'erents sur la courbe de Fargues-Fontaine} (French, with English and French summaries), Duke Math. J. \textbf{167} (2018), no. 18, 3455--3532, DOI 10.1215/00127094-2018-0034. MR3881201
\item[{[19]}] Q. Lu and W. Zheng, \emph{Duality and nearby cycles over general bases}, Duke Math. J. \textbf{168} (2019), no. 16, 3135--3213, DOI 10.1215/00127094-2019-0057. MR4027830
\item[{[20]}] T. Richarz, \emph{A new approach to the geometric Satake equivalence}, Doc. Math. \textbf{19} (2014), 209--246. MR3178249
\item[{[21]}] T. Richarz, \emph{Spaces with $\mathbb{G}_m$-action, hyperbolic localization and nearby cycles}, J. Algebraic Geom. \textbf{28} (2019), no. 2, 251--289, DOI 10.1090/jag/710. MR3912059
\item[{[22]}] P. Schneider and U. Stuhler, \emph{Representation theory and sheaves on the Bruhat-Tits building}, Inst. Hautes \'Etudes Sci. Publ. Math. \textbf{85} (1997), 97--191. MR1471867
\item[{[23]}] P. Scholze, \emph{\'Etale cohomology of diamonds}, arXiv:1709.07343, to appear in Ast\'erisque.
\item[{[24]}] P. Scholze, \emph{The local Langlands correspondence for GL$_n$ over p-adic fields}, Invent. Math. \textbf{192} (2013), no. 3, 663--715, DOI 10.1007/s00222-012-0420-5. MR3049932
\item[{[25]}] P. Scholze and J. Weinstein, \emph{Berkeley lectures on p-adic geometry}, Annals of Mathematics Studies, vol. 207, Princeton University Press, Princeton, NJ, 2020. MR4446467
\item[{[26]}] E. Viehmann, \emph{On Newton strata in the $B^+_{dR}$-Grassmannian}, Duke Math. J. \textbf{173} (2024), no. 1, 177--225, DOI 10.1215/00127094-2024-0005. MR4728690
\item[{[27]}] M.-F. Vign\'eras, \emph{Correspondance de Langlands semi-simple pour GL($n,F$) modulo $\ell\neq p$} (French), Invent. Math. \textbf{144} (2001), no. 1, 177--223, DOI 10.1007/s002220100134. MR1821157
\item[{[28]}] X. Zhu, \emph{Affine Grassmannians and the geometric Satake in mixed characteristic}, Ann. of Math. (2) \textbf{185} (2017), no. 2, 403--492, DOI 10.4007/annals.2017.185.2.2. MR3612002
\end{refs}

\medskip
\begin{center}
{\small Insitut de Math\'ematiques de Jussieu, 4 place Jussieu, 75252 Paris Cedex 05\\[0.3em]
\itshape Email address: \texttt{laurent.fargues@imj-prg.fr}}
\tn{底本此处作「Insitut」，疑为「Institut」之误。}

\medskip
{\small MPI f\"ur Mathematik, Vivatsgasse 7, 53111 Bonn, Germany\\[0.3em]
\itshape Email address: \texttt{scholze@math.uni-bonn.de}}
\end{center}
